\chapter{融合频域前馈网络的全局频域增强方法}

本章的主要贡献在于，在第三章统一骨架下提出了一种基于频域前馈网络的全局频域增强方法，以弥补局部多尺度分解在整体平均偏差优化方面的不足。与第四章侧重局部频带刻画的方案不同，本章首先将时间序列映射到复数频谱空间，在全局频域中利用可学习复数线性变换显式建模不同频率成分与变量通道之间的交互关系。其次，本文在频域映射后引入稀疏化机制，以抑制弱噪声频点并突出主要周期结构，从而提高频域表示的平滑性与鲁棒性。最后，本章设计了结合预测步长先验的门控融合机制，使全局频域特征能够按照任务需求自适应注入时域表示，为后续从实验上分析其在 MAE 和长期趋势平滑度上的优势提供了方法基础。  

\section{问题分析与研究动机}

第四章的小波包方法强调局部时频结构和多尺度分解，更适合处理局部尖峰与突变带来的大误差问题。然而，在多元长程预测任务中，除局部异常外还广泛存在较强的全局周期性和整体趋势模式。若模型主要依赖局部频带分析，虽然能够有效降低少数异常点导致的平方误差，但对整体曲线的平滑性和平均偏差控制未必最优，因此在 MAE 指标上仍然存在进一步提升空间。  

基于这一观察，本章转而关注另一类频域建模思想，即直接在全局频谱空间中学习可调节的变换。傅里叶变换能够将输入序列映射到全局频域，使不同周期成分在谱空间中得到更加明确的表达；若在此基础上进一步引入可学习的复数域线性变换与稀疏化机制，模型便可以根据数据分布自适应调整不同频率成分的强弱，并过滤对预测贡献较小的噪声频点。这样的建模方式更倾向于从整体上优化预测曲线的形状和平滑度，从而有望在 MAE 指标上取得更明显的收益。  

因此，本章在统一框架下将频域分支替换为可学习的频域前馈网络 FreTS，并重点围绕三个问题展开：一是如何在复数频谱空间中构造稳定有效的可学习变换；二是如何将稀疏频率选择与时域分支、语义分支协同融合；三是该方法是否能够在 MAE 和长期趋势预测平滑性方面相较于其他方法体现出更稳定的优势。

\section{模型架构设计}

\subsection{频域前馈网络结构设计}

本章方法以归一化后的输入序列 \(\mathbf{X}'\in\mathbb{R}^{B\times N\times L}\) 为起点，首先通过线性层将每个变量的一维时间序列投影到 \(C\) 维通道空间，得到 \([BN,L,C]\) 形式的频域前馈网络输入。随后，模型沿时间维对该表示执行实数快速傅里叶变换（RFFT），将时域序列映射为复数谱表示 \(\hat{\mathbf{X}}\in\mathbb{C}^{BN\times L_f\times C}\)，其中 \(L_f\) 为频域长度。  

在得到复数谱后，模型不直接对其进行固定规则处理，而是在频谱空间中引入可学习的复数线性变换。具体而言，FreTS 模块为实部和虚部分别设置可学习参数矩阵，通过对每个频率位置上的通道向量进行线性组合，实现对频域结构的自适应重构。与传统前馈网络在实数时域上做变换不同，这一设计将频谱中的幅值关系和相位关系一并纳入建模范围，从而使网络能够更灵活地调整不同周期模式之间的交互方式。  

经过复数域变换后，模型再对频谱表示施加非线性激活与稀疏化约束，随后通过逆傅里叶变换将结果重新映射回时域，得到频域增强后的序列表示。最后，该序列经过注意力池化与门控编码器压缩为变量级全局频域特征，并与时域分支输出共同参与后续融合。整体而言，本章方法不是将傅里叶变换作为固定的预处理步骤，而是将“频谱映射 - 可学习重构 - 回到时域”的过程纳入神经网络的可学习频域分支之中。

\subsection{复数域特征提取与门控权重融合机制}

设某一频率位置上的输入复向量为 \(\hat{\mathbf{x}}=\mathbf{x}_r+j\mathbf{x}_i\)，FreTS 模块通过两组实值参数矩阵 \(\mathbf{R}\) 和 \(\mathbf{I}\) 对其进行变换，可写为
\begin{equation}
\hat{\mathbf{y}}=(\mathbf{R}+j\mathbf{I})\hat{\mathbf{x}}.
\label{eq:chap5-complex-transform}
\end{equation}
展开后可得
\begin{equation}
\mathrm{Re}(\hat{\mathbf{y}})=\mathbf{R}\mathbf{x}_r-\mathbf{I}\mathbf{x}_i,\quad
\mathrm{Im}(\hat{\mathbf{y}})=\mathbf{R}\mathbf{x}_i+\mathbf{I}\mathbf{x}_r.
\label{eq:chap5-complex-expand}
\end{equation}
这一形式与代码中对实部和虚部分别使用参数 \texttt{r}、\texttt{i} 的实现是一致的。与简单的实数映射相比，该复数域变换不仅改变各频率成分的幅值强弱，也能够间接调整相位关系，因此更适合刻画长期周期模式之间的耦合结构。  

在复数线性层之后，本文进一步对频谱表示施加 ReLU 非线性和 SoftShrink 稀疏化操作。稀疏化的核心思想是抑制幅值较小、对预测贡献有限的频率成分，同时保留主要周期结构对应的高能量频点。阈值参数控制了稀疏程度，而初始化尺度参数则决定了复数权重在训练初期的稳定性。二者共同作用，使 FreTS 分支在保证可学习性的同时，避免频谱变换过于剧烈而破坏原始周期结构。  

本章中的“门控权重融合”主要体现在频域分支与时域分支的交互阶段。不同于第四章通过跨注意力显式引入多尺度信息，本章将时间特征、频域特征以及预测步长先验拼接后输入门控网络，生成步长感知的融合权重。这样，模型可以根据当前预测任务的跨度与上下文状态，自适应决定频域增强信息在最终表示中的注入比例。换言之，FreTS 通过稀疏化实现了频谱内部的隐式选择，又通过步长感知门控实现了分支层面的显式选择，两者共同构成了本章方法的关键设计。

\subsection{频域-时域特征交互设计}

FreTS 分支经过逆傅里叶变换后重新回到时域，其输出仍然保留与原序列一致的时间长度，因此可以被视为一种经过全局频谱重构后的增强时序表示。模型随后使用注意力池化对该序列沿时间维进行压缩，得到变量级频域特征 \(\mathbf{H}^{\text{fre}}\in\mathbb{R}^{B\times N\times C}\)，并进一步通过门控编码器强化其结构表达能力。与第四章从多个子带获得局部频带信息不同，这里得到的是更偏向全局周期和整体趋势的频域表示。  

在与时域分支融合时，本文将时域特征 \(\mathbf{H}^{\text{time}}\)、频域特征 \(\mathbf{H}^{\text{fre}}\) 以及归一化后的预测步长先验 \(\mathbf{h}^{\text{horizon}}\) 进行拼接，并通过两层感知机生成门控向量 \(\mathbf{G}^{\text{fusion}}\)。融合结果可表示为
\begin{equation}
\mathbf{H}^{\text{tf}}=\mathbf{H}^{\text{time}}+\mathbf{G}^{\text{fusion}}\odot \mathbf{H}^{\text{fre}}.
\label{eq:chap5-fusion}
\end{equation}
该设计使模型能够针对不同预测跨度动态调节频域分支的影响强度。对于较长预测步长，模型通常更依赖全局周期信息；对于较短预测步长，时域局部动态可能占主导。步长感知门控正是在这一意义上实现了频域与时域信息的任务相关融合。  

在完成时频融合后，本章方法与第三章统一骨架保持一致：冻结大语言模型生成的语义嵌入先经 \texttt{prompt\_encoder} 适配，再与时频融合表示在多头 CMA 中进行跨模态对齐；经过自适应头融合和残差插值后，解码器输出最终预测结果。在实验评价上，本章将继续汇报 MSE 与 MAE 两项指标，但重点关注 MAE 以及长期预测曲线的平滑性，以突出全局频域建模的优势。

\section{实验设置与结果分析}

\subsection{实验设置与对比基线}

本章实验设置与第四章保持一致，包括数据集划分、输入输出窗口长度、RevIN 预处理、评价指标和训练流程，以保证两章结果具备可比性。模型结构上，仅将统一框架中的频域分支替换为 FreTS 频域前馈网络，并采用默认的稀疏阈值、初始化尺度和通道配置。时域分支、语义分支、CMA 与解码器均保持不变，从而使实验结论能够更集中地反映全局频域建模本身的作用。  

为凸显本章方法的特点，对比对象建议至少包括以下几类：原始基座模型、无频域增强的统一框架、第四章的小波包版本以及必要的其他频域基线。结果表中同时给出 MSE 与 MAE，但分析时应以 MAE 为主，因为本章的核心问题在于验证可学习全局频谱建模是否能够更有效地降低平均绝对误差并改善长期预测稳定性。训练细节如优化器、学习率、batch size 和早停策略均与前文保持一致。

\subsection{MAE 指标表现及稳健性分析}

本节重点分析 FreTS 模型在 MAE 指标上的表现，并与基座模型、小波包版本及其他对比方法进行系统比较。表格设计可延续前章方式，按照“数据集 × 预测长度 × 方法”的形式汇报结果，同时保留 MSE 作为辅助指标，以便与第四章构成对照。若 FreTS 在多数数据集和较长预测步长下取得更低的 MAE，则说明全局频域建模确实更有利于降低整体平均偏差。  

在结果讨论中，应重点关注两类现象。首先，与第四章相比，FreTS 是否在 MAE 上表现出更稳定的一致性，即使在 MSE 上未必始终最优，也能在整体误差控制方面体现优势。其次，这种优势是否在周期性更明显、长期趋势更平滑的数据集上更加显著。如果这些现象得到验证，则可以较自然地引出第六章关于“全局频域建模更适合优化 MAE”的进一步分析。  

若实验包含多次重复运行，还可报告均值与标准差，用以说明 FreTS 对随机初始化的敏感程度是否较低。必要时，也可以从稀疏阈值或初始化尺度的角度补充讨论模型的稳健性，从而增强本章结论的可信度。

\subsection{长期预测趋势平滑度分析}

除数值指标外，本章还需要从预测曲线形态角度分析 FreTS 的作用。由于该方法在频谱空间中显式抑制弱噪声频点，并通过可学习变换重构全局周期信息，因此理论上更容易生成平滑、连续的长期预测曲线。为验证这一点，可在较长预测步长设置下选取代表性样本，绘制真实值、FreTS 预测值以及若干基线方法预测值的对比曲线，并观察各方法在趋势跟随性、局部抖动程度和整体平滑性上的差异。  

若需要给出更定量的分析，还可以引入一阶差分或二阶差分的范数作为平滑度指标，用于衡量预测序列在相邻时间步之间的波动程度。若 FreTS 模型在保持较好拟合的前提下同时表现出更低的差分波动，则可进一步说明其确实通过全局频域建模和稀疏化机制减少了不必要的高频噪声。这样的分析不仅有助于解释 MAE 改善的来源，也能增强本章关于长期趋势预测能力的论证力度。

\subsection{消融实验：稀疏阈值与初始化尺度的影响}

由于本文实现中的 FreTS 频域前馈网络采用单层核心频谱映射结构，因此本节将消融重点放在两个更贴近实际实现的关键超参数上，即稀疏阈值 \texttt{sparsity\_threshold} 和初始化尺度 \texttt{scale}。其中，前者决定了频谱稀疏化的强弱，后者影响复数线性权重在训练初期的分布范围。通过分别调整这两个参数并观察 MAE、MSE 以及训练稳定性的变化，可以更准确地分析 FreTS 模块为何能够在当前配置下取得较优结果。  

从理论上看，过小的稀疏阈值可能无法有效抑制噪声频点，使频域分支保留过多无效成分；而过大的阈值又可能误删有用的周期信息，导致表示能力下降。同样，初始化尺度过小可能限制模型早期的表达能力，过大则可能导致频谱变换不稳定。因此，本节消融的目标是找到一组兼顾稳定性与表达性的参数配置。若实验结果表明当前默认设置在多数数据集上取得较优的 MAE，并且训练过程更平稳，则可进一步支撑本章设计的合理性。

\section{本章小结}

本章在统一框架下提出了一种基于 FreTS 的全局频域增强方法。该方法通过将输入序列映射到复数频谱空间，在频域中引入可学习线性变换和稀疏化机制，再结合步长感知门控实现与时域分支的自适应融合，从而构建起一种兼顾全局周期建模和任务相关选择性的频域建模路径。  

相较于第四章侧重局部多尺度分解的方案，本章方法更强调从整体频谱结构中提取对长期预测有用的稳定模式。若实验结果表明 FreTS 在 MAE 和长期趋势平滑度方面具有更明显优势，则可以据此说明：可学习的全局频域重构对于降低整体平均误差具有独特价值，并为第六章进一步解释“小波包偏向优化 MSE、频域前馈网络偏向优化 MAE”的现象提供了第二条关键证据链。
